\subsubsection{FingerNet}
\label{sec:low_level_finger}
In this stage, we generate finger motions for the pre-contact and post-contact phases to create natural transitions between the rest and grasp poses, forming a complete interaction sequence along with predictions in prior stages.
Given the start and end finger poses and the wrist trajectory in between, we employ a conditional diffusion model to predict the finger motions denoted as $\bm{F} \in \mathbb{R}^{T\times D'}$, representing local 6D rotations~\cite{zhou2019continuity}.

% Similar to ManipNet~\cite{zhang2021manipnet}, we only process the finger motions for the right hand. 
% To handle the left hand, we mirror the entire scene, process the mirrored left hand as if it were the right hand, and then mirror the output back.
% This approach helps counteract the limited scale of the dataset~\cite{GRAB:2020}.
% At training time, both hand data are given so there is no bias on handedness.

% \paragraph{Motion Representation} 
% Finger motion $\boldsymbol{F} \in \mathbb{R}^{T\times D'}$ contains local joint rotations in 6D representation. 

% \paragraph{Hand-Object Spatial Representation.} 


 
% \paragraph{Input Condition Representation.} 
The model is conditioned on $\bm{c}_{f} = \left\{\bm{P}, \bm{F}_{s}, \bm{F}_{e}\right\}$, where $\bm{P}$ captures spatial hand-object relationships and $\bm{F}_{s}, \bm{F}_{e}$ are the start and end finger poses. Following prior work~\cite{zhang2021manipnet,taheri2023grip}, $\bm{P} \in \mathbb{R}^{T \times 100}$ is computed by sampling 100 palm-side mesh vertices and measuring their closest distances to the object, with the finger pose set to mean pose when calculating $\bm{P}$. 
% \jiaman{Is the data representation P the same as prior work? We probably need to clearly say what is prior work already have and what is new thing we proposed. So is P the same, but conditional diffusion model with P is new?} \zhen{Yes P is the same, but conditional diffusion model with P is new. Fixed.}


% \paragraph{Mirroring Strategy.} 
% The FingerNet is trained to process only the right hand. To generate finger motions for the left hand, we apply a mirroring strategy similar to~\cite{zhang2021manipnet}.
% Specifically, we first mirror the left hand and the corresponding object mesh to simulate a right-hand scenario. We then process this mirrored data to generate the finger motions. Finally, we mirror the resulting motions back to obtain the finger movements of the left hand. 

% To obtain the training data, we extracted all hand-reaching and hand-leaving motions from the hand-object interaction dataset. 
% The model is then optimized with the "simple" objective \cite{ho2020denoising} as described in Equation \ref{equ:simple}.

% During inference, we utilize the wrist trajectory generated by the interaction module, along with the resting finger pose and the precise finger pose $\hat{\boldsymbol{\theta}}$, as conditions for generating finger movements.








